% G3-4 Pembagian data berurutan waktu dan purging.
% Sumber: diolah penulis.
\begin{tikzpicture}[x=1mm,y=-1mm,font=\footnotesize,
    latih/.style={draw,minimum height=8mm,inner sep=0pt,anchor=west},
    buang/.style={draw,pattern=north east lines,minimum height=8mm,
      inner sep=0pt,anchor=west},
    uji/.style={draw,fill=black!18,minimum height=8mm,inner sep=0pt,
      anchor=west},
    panah/.style={-{Stealth[length=2.2mm]}}]

  % --- arah waktu ---
  \draw[panah] (14,-16) -- (124,-16);
  \node[anchor=south,font=\footnotesize] at (69,-16)
    {arah waktu (tidak ada pengacakan)};

  % ===== (a) pembagian utama =====
  % Batas periode pengembangan ada di x = 88 untuk semua baris.
  \node[anchor=south west] at (14,-6) {\textbf{(a)} pembagian utama};
  \node[latih,minimum width=74mm] (a1) at (14,6) {periode pengembangan (80\%)};
  \node[buang,minimum width=6mm]  (a2) at (88,6) {};
  \node[uji,minimum width=30mm]   (a3) at (94,6) {blok uji final (20\%)};

  % ===== (b) panjang data latih pada Tahap ketiga =====
  \node[anchor=south west] at (14,22)
    {\textbf{(b)} panjang data latih pada Tahap ketiga};

  \node[anchor=east,font=\small] at (12,30) {100\%};
  \node[latih,minimum width=74mm] (b1) at (14,30) {data latih};
  \node[buang,minimum width=6mm]       at (88,30) {};
  \node[uji,minimum width=30mm]        at (94,30) {blok uji final};

  \node[anchor=east,font=\small] at (12,44) {75\%};
  \node[latih,minimum width=55.5mm] (c1) at (32.5,44) {data latih};
  \node[buang,minimum width=6mm]         at (88,44)   {};
  \node[uji,minimum width=30mm]          at (94,44)   {blok uji final};

  \node[anchor=east,font=\small] at (12,58) {50\%};
  \node[latih,minimum width=37mm] (d1) at (51,58) {data latih};
  \node[buang,minimum width=6mm]  (d2) at (88,58) {};
  \node[uji,minimum width=30mm]        at (94,58) {blok uji final};

  % --- label purging, menunjuk ke potongan yang diarsir ---
  \draw[panah] (72,76) -- (d2.south);
  \node[anchor=north east,align=right] at (71,74) {\textit{purging}};

  % --- legenda ---
  \node[latih,minimum width=6mm] at (14,88) {};
  \node[anchor=west] at (21,88) {periode pengembangan};
  \node[buang,minimum width=6mm] at (58,88) {};
  \node[anchor=west] at (65,88) {dibuang (\textit{purging})};
  \node[uji,minimum width=6mm] at (100,88) {};
  \node[anchor=west] at (107,88) {blok uji final};

\end{tikzpicture}
